\documentclass[12pt, a4paper]{ctexart}
\usepackage{fontspec}
\usepackage{xcolor}
\usepackage{hyperref}
\usepackage{geometry}
\usepackage{enumitem}
\usepackage{parskip}
\usepackage{titlesec}
\usepackage{tcolorbox}
\usepackage{setspace}
\usepackage{listings}
\usepackage{amssymb}

\geometry{left=2.5cm, right=2.5cm, top=2.5cm, bottom=2.5cm}

\setmainfont{Times New Roman}
\setCJKmainfont{SimSun}

\hypersetup{
    colorlinks=true,
    linkcolor=blue!70!black,
    urlcolor=cyan,
    pdftitle={Java开发入门-逐题精讲解义},
    pdfauthor={专升本-fifine老师}
}

\definecolor{myblue}{RGB}{30,100,200}
\definecolor{lightblue}{RGB}{230,240,255}
\definecolor{darkblue}{RGB}{0,50,120}

\newtcolorbox{bluebox}[1][]{
    colback=lightblue,
    colframe=myblue,
    coltitle=darkblue,
    fonttitle=\bfseries,
    title={#1},
    boxrule=0.8pt,
    arc=4pt,
    left=6pt,
    right=6pt,
    top=6pt,
    bottom=6pt,
    before skip=8pt,
    after skip=8pt
}

\titleformat{\section}
  {\normalfont\Large\bfseries\color{myblue}}{\thesection}{1em}{}
\titlespacing*{\section}{0pt}{3.5ex plus 1ex minus .2ex}{1.5ex plus .2ex}

\title{\textcolor{myblue}{\textbf{Java开发入门 - 逐题精讲解义}}}
\author{专升本-fifine老师 教学组}
\date{2026年03月05日}

\lstset{
    language=Java,
    basicstyle=\ttfamily\small,
    keywordstyle=\color{blue},
    commentstyle=\color{green!60!black},
    stringstyle=\color{orange},
    numbers=left,
    numberstyle=\tiny\color{gray},
    frame=single,
    breaklines=true,
    columns=fullflexible,
    showstringspaces=false
}

\newcommand{\code}[1]{\texttt{#1}}

\begin{document}

\maketitle
\thispagestyle{empty}
\newpage

\section{Java开发入门}

\begin{enumerate}[label=\textbf{\arabic*.}]

	\item \textbf{【本章说明】}

	      \begin{bluebox}{章节概述}
		      \textbf{题目数量}：暂无本章题目

		      \textbf{章节内容}：

		      本章主要考察Java环境搭建与运行机制，包括：
		      \begin{itemize}[label={}, leftmargin=*, nosep]
			      \item JDK/JRE/JVM的区别与联系
			      \item Java环境的安装与配置
			            \textbackslashitem Java程序的编译与运行过程
			      \item \texttt{javac}与\texttt{java}命令的使用
			      \item 类文件的结构与字节码
		      \end{itemize}

		      \textbf{命题人视角}：

		      虽然当前题库暂无本章节题目，但本章是Java学习的基础，也是后续章节的前置知识。可能的考点方向包括：
		      \begin{itemize}[label={}, leftmargin=*, nosep]
			      \item 考查对JVM内存模型的基本理解
			      \item 考查Java程序编译与运行的完整流程
			      \item 考查类加载机制的基本概念
		      \end{itemize}

		      \textbf{【应背诵知识点】}

		      \begin{enumerate}[label={}, leftmargin=*, nosep]
			      \item \textbf{JDK/JRE/JVM三层架构}：
			            \begin{itemize}
				            \item JDK（Java Development Kit）：开发工具包，包含JRE + 开发工具（javac、javadoc等）
				            \item JRE（Java Runtime Environment）：运行环境，包含JVM + 核心类库
				            \item JVM（Java Virtual Machine）：虚拟机，运行Java字节码
			            \end{itemize}

			      \item \textbf{Java程序运行过程}：
			            \begin{itemize}
				            \item 编译：\texttt{javac Xxx.java} → \texttt{Xxx.class}
				            \item 运行：\texttt{java Xxx}
				            \item JVM加载.class文件，解释或编译为机器码执行
			            \end{itemize}

			      \item \textbf{关键字考点}：
			            \begin{itemize}
				            \item \texttt{public}：访问修饰符，表示公开访问
				            \item \texttt{class}：定义类的关键字
				            \item \texttt{static}：静态修饰符，属于类而非对象
				            \item \texttt{void}：返回类型，表示无返回值
				            \item \texttt{main}：程序入口方法名
				            \item \texttt{String[] args}：命令行参数数组
			            \end{itemize}
		      \end{enumerate}
	      \end{bluebox}

\end{enumerate}

\end{document}
